\chapter{发送通路}\label{ch:tx}

本章按数据流向说明发送核心 \code{afdx_tx_core} 的五个部分：分页缓存 \code{afdx_tx_buf} 与空闲页池 \code{afdx_tx_free}、调度引擎 \code{afdx_tx_engine}、组帧器 \code{afdx_tx_builder} 和双网 MAC \code{afdx_mac_tx}。

\section{发送核心与配置表}\label{sec:txtables}

图~\ref{fig:tx} 是发送核心的框图。应用和服务的消息经 \code{afdx_tx_buf} 分页写入载荷 RAM 并登记入队；\code{afdx_tx_engine} 按 BAG 选出下一帧并授权给 \code{afdx_tx_builder}；组帧器预读载荷、拼接帧头，逐字节交给 \code{afdx_mac_tx} 发往两网。

\figN[发送核心框图]{f07_tx}{发送核心及其配置、状态、载荷存储之间的连接}{fig:tx}

发送侧共 7 张表，全部由 \code{afdx_ram_sdp}（一写一读、同步读）例化，每张表的写口和读口各只归一个模块，见表~\ref{tab:txtables}。

\begin{table}\centering\small
\caption{发送侧配置与状态表}\label{tab:txtables}
\begin{tabularx}{\linewidth}{>{\raggedright\arraybackslash}p{2.1cm}>{\raggedleft\arraybackslash}p{1.9cm}>{\raggedright\arraybackslash}p{1.9cm}>{\raggedright\arraybackslash}p{1.9cm}L}
\toprule
表 & 深度$\times$位宽 & 写口 & 读口 & 内容 \\
\midrule
\code{tx_cfg} & 1024$\times$32 & 配置总线 & tx\_engine & 每条 VL 8 个字，地址 \{VL[6:0], 字号[2:0]\} \\
\code{tx_state} & 128$\times$63 & tx\_engine & tx\_engine & \code{timer_active}、\code{next_time}、\code{sn}、\code{q_count}、\code{q_head}、\code{q_tail} \\
\code{tx_mnext} & 256$\times$8 & tx\_engine & tx\_engine & 同一 VL 队列中的下一条消息（首页号） \\
\code{tx_desc} & 256$\times$44 & tx\_buf & tx\_engine & 消息描述：长度、\code{port_en}、源/目的端口 \\
\code{tx_pnext} & 256$\times$8 & tx\_buf & tx\_builder & 消息内的下一页 \\
\code{tx_payload} & 131072$\times$8 & tx\_buf & tx\_builder & 128~KiB 载荷，地址 \{页号[7:0], 偏移[8:0]\} \\
\code{tx_free} & 256$\times$8 & tx\_free & tx\_free & 空闲页号 FIFO \\
\bottomrule
\end{tabularx}
\end{table}

\code{tx_cfg} 中每条 VL 占 8 个字，其中 6 个配置字有效，字段见表~\ref{tab:cfgfmt}。配置字合并在同一张 1024$\times$32 的表中，相比拆分为六张浅表减少了块 RAM 占用；引擎在授权操作中逐拍读取六个配置字，步骤见表~\ref{tab:txops}。

\code{afdx_tx_core} 例化 7 张表和 5 个子模块，并持有 \code{tx_cfg} 的写口：复位后先用 1024 拍把 \code{tx_cfg} 清零，此后只在 \code{cfg_we} 与引擎的 \code{cfg_ok} 同拍为 1 时寄存写地址与数据，并在下一拍写入配置表。

\code{tx_cfg} 清零、空闲页装满和 \code{tx_state} 清零三者都完成后 \code{init_done} 置 1，引擎和分页缓存才开始工作。

\section{空闲页池}\label{sec:txfree}

\code{afdx_tx_free} 管理载荷缓存的 256 个页号；每页从分配到归还期间由一条消息独占。它是一块 256$\times$8 的 RAM 加读写指针和计数组成的 FIFO，复位后用 257 拍依次写入页号 0--255，装完后开放取页。

页池提供一个取页口和两个还页口，均带握手：
\begin{itemize}
  \item 取页口 \code{pop_req}/\code{pop_ready}/\code{pop_ack}/\code{pop_ok}/\code{pop_page} 接 \code{tx_buf}，请求应答式；
  \item 还页口 1 \code{ret1_*} 接 \code{tx_builder}，读完一页还一页；
  \item 还页口 2 \code{ret2_*} 接 \code{tx_buf}，放弃消息时退页。
\end{itemize}

取页由四步状态机完成，见表~\ref{tab:freefsm}。从请求被接受到 \code{pop_ack} 共 4 拍，其中 2 拍用于同步 RAM 读出和寄存。

\begin{table}\centering\small
\caption{空闲页池取页状态}\label{tab:freefsm}
\begin{tabularx}{\linewidth}{>{\raggedright\arraybackslash}p{2.4cm}L>{\raggedright\arraybackslash}p{2.6cm}}
\toprule
状态 & 动作 & 转移 \\
\midrule
\code{POP_IDLE} & \code{pop_ready}$=1$；接受请求时若 \code{count}$=0$ 立即应答 \code{pop_ok}$=0$，否则发出读 \code{rd_ptr} & 有页 $\rightarrow$ \code{POP_READ} \\
\code{POP_READ} & 等待同步 RAM 读出 & $\rightarrow$ \code{POP_CAPTURE} \\
\code{POP_CAPTURE} & 寄存读出的页号 & $\rightarrow$ \code{POP_REPLY} \\
\code{POP_REPLY} & 给出 \code{pop_ack}、\code{pop_ok}$=1$ 和页号 & $\rightarrow$ \code{POP_IDLE} \\
\bottomrule
\end{tabularx}
\end{table}

每个还页口先把页号锁存至各自的保持寄存器（\code{ret*_ready} 为该寄存器空），再经共用写口逐个写入 FIFO 尾部；两个寄存器均满时还页口 1 优先。取页与还页同拍发生时 \code{count} 不变。

FIFO 深度等于页数，任一页号在池中至多出现一次。写口每拍可写入一个页号，两个还页口各自最快每 2 拍还一页。

\section{分页缓存}\label{sec:txbuf}

\code{afdx_tx_buf} 从消息接口逐字节收下消息，按页写入 \code{tx_payload}，在 \code{tx_pnext} 中把同一消息的页串成链，最后写 \code{tx_desc} 并向引擎入队。载荷缓存共 256 页、每页 512 字节，地址为 \{页号[7:0], 偏移[8:0]\}。

分页缓存的接口包括消息接口 \code{m_*}（表~\ref{tab:ifm}）；取页口和还页口 2；\code{tx_payload}、\code{tx_pnext}、\code{tx_desc} 的写口；入队请求应答口（表~\ref{tab:ifgrant}）；五个统计脉冲 \code{ev_accept}、\code{ev_drop_dis}、\code{ev_drop_full}、\code{ev_drop_long}、\code{ev_drop_nopage}。

\subsection{消息接收与页分配}

图~\ref{fig:txbuf} 是其状态机，各状态的动作见表~\ref{tab:txbuffsm}。\code{m_ready} 由状态组合产生，初始化完成后只在 IDLE、RECV、DRAIN 三个状态为 1；字节传输以 \code{m_valid} 与 \code{m_ready} 同时为 1 为条件。

\figN[分页缓存状态机]{f19_txbuf}{分页缓存的消息接收、页分配、入队与异常退页流程}{fig:txbuf}

\begin{table}\centering\small
\caption{分页缓存各状态}\label{tab:txbuffsm}
\begin{tabularx}{\linewidth}{>{\raggedright\arraybackslash}p{2.1cm}L>{\raggedright\arraybackslash}p{4.0cm}}
\toprule
状态 & 动作 & 转移 \\
\midrule
IDLE & 收下首字节存入 \code{held_data}，取样 VL 与端口，长度记 1 & $\rightarrow$ ALLOC \\
ALLOC & 置 \code{pop_req} & $\rightarrow$ WAIT\_PAGE \\
WAIT\_PAGE & 等 \code{pop_ack}；得页后记入 6 个页寄存器之一，首页记为 \code{enq_head}，否则写 \code{tx_pnext}[上一页] $=$ 新页；把暂存字节写到新页偏移 0 & 无页 $\rightarrow$ RETURN；暂存字节是末字节 $\rightarrow$ DESC；否则 $\rightarrow$ RECV \\
RECV & 每拍收一字节写入当前页，偏移加 1；偏移到 511 时置 \code{boundary} & 已满 1471 字节仍有数据 $\rightarrow$ RETURN；跨页字节暂存 $\rightarrow$ ALLOC；末字节 $\rightarrow$ DESC \\
DESC & 写 \code{tx_desc}[\code{enq_head}] $=$ \{长度, \code{port_en}, 端口\} & $\rightarrow$ ENQUEUE \\
ENQUEUE & 置 \code{enq_valid} & $\rightarrow$ REPLY \\
REPLY & 等 \code{enq_ack}，记下 \code{enq_result} & 收下 $\rightarrow$ EVENT；拒绝 $\rightarrow$ RETURN \\
RETURN & 经还页口 2 逐页退回已用页 & 消息已结束 $\rightarrow$ EVENT；否则 $\rightarrow$ DRAIN \\
DRAIN & 继续收字节但不保存 & 末字节 $\rightarrow$ EVENT \\
EVENT & 按原因给出一个统计脉冲 & $\rightarrow$ IDLE \\
\bottomrule
\end{tabularx}
\end{table}

写满一页后，将下一字节暂存于 \code{held_data}，取到新页后写入其偏移 0。取页由跨页字节到达触发。最长消息占 3 页，6 个页寄存器为更小的页长参数留有余量。

\FloatBarrier
\subsection{时序与异常处理}

以首字节被收下为第 0 拍：第 1 拍发取页请求，第 2 拍被空闲页池接受，第 6 拍 \code{pop_ack} 到达并写入首字节，第 7 拍回到 RECV。每取一页 \code{m_ready} 为 0 共 6 拍。消息末字节之后，DESC、ENQUEUE 各 1 拍，引擎入队 6 拍后应答，EVENT 1 拍，\code{m_ready} 共为 0 约 10 拍。

超长（第 1472 个字节到来）、取页失败与入队被拒三种情况均转入 RETURN 态：先退回本消息已用的全部页，再继续接收并丢弃剩余字节。被丢弃消息占用的页在其统计脉冲之前全部退回。\code{reason} 记录首次出现的原因，每条消息在 EVENT 状态产生一个统计脉冲。

每条消息的统计脉冲出现在末字节被收下后 96 拍以内、下一条消息末字节之前，因此第 $k$ 个脉冲对应第 $k$ 条消息。

\code{m_ready} 连续为 0 的上界由消息末尾决定：$2 + 64 + 9 + 1 = 76$ 拍，取设计约束 96 拍（第~\ref{sec:budget}~节）。一条 1471 字节消息在应用口共占 1499 拍，短于它在线路上的 1538 拍。

\section{调度引擎}\label{sec:txeng}

\code{afdx_tx_engine} 是 \code{tx_state} 和 \code{tx_mnext} 的唯一读写控制模块，负责入队判决、按 BAG 选帧、授权组帧器和根据实际起始时刻更新下一次允许时刻。

\subsection{队列状态与操作仲裁}

\code{tx_state} 每条 VL 一项 63 位，见表~\ref{tab:txstate}。每条 VL 的消息队列是一条单向链表：\code{q_head}、\code{q_tail} 存首尾消息的首页号，\code{tx_mnext}[首页号] 指向下一条消息，因此 128 条 VL 共用 256 项链表。

\begin{table}\centering\small
\caption{\texttt{tx\_state} 字段}\label{tab:txstate}
\begin{tabularx}{\linewidth}{>{\centering\arraybackslash}p{1.6cm}>{\raggedright\arraybackslash}p{2.4cm}L}
\toprule
位 & 字段 & 含义 \\
\midrule
62 & \code{timer_active} & \code{next_time} 有效；为 0 表示可立即发送 \\
61:30 & \code{next_time} & 下一帧最早起始的 \code{now} 值 \\
29:22 & \code{sn} & 下一帧的序号，复位为 0，之后 1--255 循环 \\
21:16 & \code{q_count} & 排队消息数，不超过配置的 \code{qmax}（至多 63） \\
15:8 & \code{q_head} & 队首消息首页号 \\
7:0 & \code{q_tail} & 队尾消息首页号 \\
\bottomrule
\end{tabularx}
\end{table}

引擎一次只做一种操作，见图~\ref{fig:txfsm}。IDLE 下按“起始确认 $>$ 入队 $>$ 授权 $>$ 扫描”的优先级选择，各操作按 \code{phase} 逐拍推进，做完即回 IDLE，见表~\ref{tab:txops}。

\code{want_enq} 表示入队请求有效且未处理，并满足已有候选或扫描配额已满；\code{want_grant} 表示已有候选、授权端就绪且无待起始帧。无上述操作时，IDLE 每 2 拍检查一条 VL。

\figN[发送调度引擎状态机]{f08_txfsm}{发送调度引擎的起始确认、消息入队、帧授权与扫描流程}{fig:txfsm}

\begin{table}\centering\small
\caption{调度引擎的操作}\label{tab:txops}
\begin{tabularx}{\linewidth}{>{\raggedright\arraybackslash}p{1.5cm}>{\raggedright\arraybackslash}p{2.4cm}L>{\raggedleft\arraybackslash}p{1.1cm}}
\toprule
操作 & 触发 & 逐拍动作 & 拍数 \\
\midrule
START & 组帧器报告 $t_\text{start}$（先锁存到 \code{ack_pending}） & 0：读 \code{tx_state}[\code{pend_vl}]；1：保存；2：写 \code{timer_active}$=1$、\code{next_time}$=t_\text{start}+$\code{bag}，清 \code{pend_valid} & 3 \\
ENQ & \code{enq_valid}，且已有候选或扫描配额已满 & 0：读 \code{tx_cfg} 字 0 和 \code{tx_state}；1：保存；2：$d=$\code{now}$-$\code{next_time}；3：$d+512$；4：判决，收下则写回 \code{q_count}$+1$、新队尾，非空队列另写 \code{tx_mnext}[旧队尾]，并在条件满足时直接定为候选 & 5 \\
GRANT & 有候选且 \code{grant_ready} & 0--4：依次读配置字 1--5；1：保存字 0 和 \code{tx_state}；2：读 \code{tx_desc}、\code{tx_mnext}[队首]；4：保存描述；7：给出授权，写回 SN 前进、\code{q_count}$-1$、队首前移，记 \code{pend_vl} & 8 \\
扫描 & 无上述请求且无候选 & 每 2 拍发出一条 VL 的读，经 5 级流水判定（图~\ref{fig:scan}） & 2/VL \\
\bottomrule
\end{tabularx}
\end{table}

\FloatBarrier
\subsection{入队与扫描}

入队判决按顺序检查：VL 未启用或两网都未启用时返回结果码 1；$n$ 大于 \code{max_payload} 时返回 3；\code{q_count} 不小于 \code{qmax} 时返回 2；其余情况接受入队并返回 0。收下时若当前没有候选、该 VL 不是 \code{pend_vl} 且已在 512 拍窗口内，直接定为候选，省去一次扫描。

扫描每 2 拍发出一条 VL 的 \code{tx_state} 读，结果经图~\ref{fig:scan} 的 5 级流水判定。记 $d=$\code{now}$-$\code{next_time}，判据为：
\begin{itemize}
  \item \code{q_count}$\ne 0$、（\code{timer_active}$=0$ 或 $d+$\code{LOOKAHEAD}$\ge 0$）且不是 \code{pend_vl}：定为候选，扫描停在下一条；
  \item \code{q_count}$=0$、\code{timer_active}$=1$ 且 $d\ge 0$：写回 \code{timer_active}$=0$，防止 \code{now} 回绕约 34~s 后把旧时刻误判为将来；
  \item 判定完成时 \code{scan_commit} 指针前移至下一条 VL；任一操作插入时流水线作废，自 \code{scan_commit} 重新发出，并按更新后的状态重新判定。
\end{itemize}

\figN[扫描流水]{f21_scan}{扫描流水。每 2 拍发出一条 VL，第 $c+5$ 拍给出判定；START、ENQ 或 GRANT 插入时清空流水，从 \texttt{scan\_commit} 重新扫描}{fig:scan}

扫描从上次选中的下一条 VL 继续，128 条 VL 轮流被检查。入队操作仅在已有候选或扫描满 16 条（\code{SCAN_QUOTA}）时插入，使持续入队期间仍能推进扫描；无候选时，任一 VL 相邻两次被检查的间隔约 400 拍（第~\ref{sec:budget}~节）。

\subsection{授权状态与发送时刻}

一帧从授权到真正起始之间，该 VL 的 \code{next_time} 仍是旧值。\code{pend_vl} 保存该帧的 VL，扫描和入队将其排除在候选之外；起始确认写回新的 \code{next_time} 时清除该标记。组帧器的一个“下一帧”位置对应一个 \code{pend_vl} 寄存器。

128 条 VL 扫描一遍 256 拍，长于最短帧的 84 拍时隙；若到期后才开始找，发送时刻会随扫描位置浮动。引擎在到期前 \code{LOOKAHEAD}$=512$ 拍内选出候选并完成授权，组帧器备好长度和校验和，\code{now} 到达允许时刻 $T$ 时 MAC 恰好发出第一个前导字节。回访约 400 拍加授权和准备 22 拍小于 512，见第~\ref{sec:budget}~节和图~\ref{fig:pipe}。

\code{next_time} 以实际起始时刻 $t_\text{start}$ 加 \code{bag} 计算。线路争用使某帧晚发时，下一帧的允许时刻随之顺延，同一 VL 相邻两帧的起始间隔始终不小于 BAG。

授权时把当前 \code{sn} 交给组帧器并写回下一值：复位后每条 VL 第一帧为 0，其后 1、2、…、255、1。

\section{组帧器}\label{sec:txbld}

\code{afdx_tx_builder} 把授权的字段组成完整帧内容（目的 MAC 到 SN），在允许时刻交给 MAC，并把实际起始时刻回报给引擎。

它内部最多同时有两帧：一帧在 \code{stream_active} 状态向 MAC 送字节，另一帧在 \code{pending} 位置准备并等待。\code{grant_ready}$=\neg$\code{pending}，\code{pending} 在本帧的 \code{tx_sof} 那拍清除，因此下一帧的授权可在当前帧起始后立即进入，准备在当前帧发送期间完成。

\subsection{帧头准备与载荷发送}

组帧器的结构见图~\ref{fig:builder}：帧头走 336 位移位寄存器，载荷走预读流水和 8 项字节 FIFO，补零和 SN 由计数器产生，三者经一个选择器送出。

\figN[组帧器数据通路]{f22_builder}{组帧器数据通路。帧头与载荷分两路，在选择器汇合后交给 MAC；页归还由独立的读元数据流水触发，与载荷进入字节 FIFO 同时发生}{fig:builder}

收下授权的同拍锁存全部字段和 \code{USER_ID}，IP 标识取 16 位计数器并加 1。之后 10 拍算 IP 头校验和，每拍只做一次加法，见表~\ref{tab:prep}。

\begin{table}\centering\small
\caption{组帧器准备步骤}\label{tab:prep}
\begin{tabularx}{\linewidth}{>{\centering\arraybackslash}p{1.4cm}L}
\toprule
\code{prep} & 动作 \\
\midrule
0 & IP 总长 $=n+28$，UDP 长 $=n+8$；累加器置初值 \code{0x4500}$+$\code{0x0111}$=$\code{0x04611} \\
1--6 & 依次累加 IP 总长、标识、源 IP 高/低 16 位、目的 IP 高/低 16 位 \\
7--8 & 两次把高 4 位进位加回低 16 位 \\
9 & 取反得校验和，置 \code{prepared} \\
\bottomrule
\end{tabularx}
\end{table}

起始条件为 \code{prepared}、上一帧已送完，且 \code{grant_timed}$=0$、允许时刻已过或距允许时刻至多 4 拍。4 拍提前量补偿起始路径上的四级寄存器：\code{now} 寄存 1 拍、差值 \code{now_q}$-T$ 寄存 1 拍、\code{s_valid} 寄存 1 拍、MAC 判断起始并寄存输出 1 拍。于是线路空闲时第一个前导字节恰在 \code{now}$=T$ 那一拍出现。

起始拍把 42 字节帧头并行装入移位寄存器，A、B 两路共用，只在第 11 字节（接口号）B 网换为 \code{40}。

MAC 每取一字节，按 \code{part} 依次送出帧头 42 字节、载荷 $n$ 字节、$n<17$ 时补零到 17 字节、带 \code{s_last} 的 SN。\code{tx_sof} 那拍的 \code{now} 作为 $t_\text{start}$ 回报引擎。

起始拍起，只要 FIFO 中不足 5 字节就每拍发出一个 \code{tx_payload} 读；读请求寄存、同步 RAM 和输出寄存共 3 拍后写入 FIFO。

当前页与下一页号分开保存：起始时读 \code{tx_pnext}[首页]，每次跨页再读下一页的后继，跨页时请求连续不断。读到页末字节或消息末字节时，该页号随数据一起进入 4 项归还 FIFO，经还页口 1 退回空闲页池。

\subsection{时序与异常处理}
\begin{itemize}
  \item \textbf{FIFO 容量：}发出读时 FIFO 中至多 4 字节，另有至多 3 个读在途，最多 $4+3+1=8$ 字节，等于 FIFO 深度。
  \item \textbf{供数时序：}MAC 在起始后第 8 拍才取第 0 个字节，帧头又占 42 拍，第一个载荷字节在起始后第 50 拍才被取走，而 FIFO 在起始后第 4 拍就开始填充；此后每拍补入一字节、取出一字节。下一页号在跨页前 500 拍以上已读出。
  \item \textbf{准备时间：}下一帧准备需 11 拍，可与当前帧的发送重叠；最短帧时隙为 84 拍。
\end{itemize}
从消息末字节被收下到 \code{tx_sof} 的名义值为 31 拍，上界 100 拍，逐拍分解见第~\ref{sec:budget}~节图~\ref{fig:pipe}。

MAC 报 \code{tx_underrun} 时本帧记为坏帧，不产生 \code{ev_tx_a}/\code{ev_tx_b}；其余帧在 \code{tx_eof} 时按本帧启用的网络各给一个脉冲。

\Needspace{14\baselineskip}
\section{双网发送 MAC}\label{sec:mactx}

\code{afdx_mac_tx} 用一个状态机驱动 A、B 两套 CRC 和输出寄存器，两网同拍输出。发送节拍由 MAC 控制，上游按其握手信号提供字节。图~\ref{fig:macfsm}(a) 是其状态机，(b) 是第~\ref{sec:macrx}~节的接收 MAC。

eLinx 生成的三速以太网 IP 带用户侧 FIFO、保留 FCS，且每网一个实例；相同器件下的资源对照见测试报告。

发送状态机仅在 \code{byte_en} 有效时推进。接收状态机使用寄存后的有效标记和数据；PREAMBLE 或 \code{BAD_PRE} 因有效标记撤销而退出时，报告前导错误。

各状态的动作见表~\ref{tab:macfsm}，发送一帧的字节时间见表~\ref{tab:macbyte}，$L$ 为取走的字节数。起始那一拍取样 \code{s_en_a}、\code{s_en_b}；未启用的网络在整帧期间 \code{tx_valid} 与 \code{tx_data} 保持为 0；两网均未启用时状态机仍按相同节拍推进，使各帧时序与网络使能无关。上游一直有帧等候时，帧间隔恰为 12 拍。

数据阶段若上游未备好字节（欠载），MAC 给出 \code{tx_underrun} 脉冲，发出 4 字节不取反的错误 FCS，然后继续接收并丢弃上游这一帧的剩余字节。

复位期间线路输出全为 0，复位结束后先等满 12 拍帧间隔再允许首帧。

\figN[双网 MAC 状态机]{f10_macfsm}{双网 MAC 状态机：(a) 发送及欠载处理；(b) 接收及前导检查}{fig:macfsm}

\begin{table}\centering\small
\caption{发送 MAC 各状态}\label{tab:macfsm}
\begin{tabularx}{\linewidth}{>{\raggedright\arraybackslash}p{2.1cm}L>{\raggedright\arraybackslash}p{3.4cm}}
\toprule
状态 & 动作 & 转移 \\
\midrule
GAP & 输出 0；\code{gap_left} 递减 & \code{gap_left}$=0$ 且 \code{s_valid} $\rightarrow$ PREAMBLE，同拍取样两网使能、给 \code{tx_sof} \\
PREAMBLE & 输出 7 个 \code{55} 和 \code{D5}；末拍 \code{s_ready}$=1$ 取走第 0 字节 & 取到字节 $\rightarrow$ DATA；未备好 $\rightarrow$ FCS（欠载） \\
DATA & 每拍输出一字节并更新两网 CRC & 末字节已出 $\rightarrow$ FCS；未备好 $\rightarrow$ FCS（欠载） \\
FCS & 输出 4 字节 FCS，低字节先发；第 4 拍后撤 \code{valid}、给 \code{tx_eof}，\code{gap_left}$=11$ & 正常 $\rightarrow$ GAP；欠载 $\rightarrow$ DRAIN \\
DRAIN & \code{s_ready}$=1$，继续接收并丢弃上游剩余字节 & 收到 \code{s_last} $\rightarrow$ GAP \\
\bottomrule
\end{tabularx}
\end{table}

\begin{table}\centering\small
\caption{MAC 发送的字节时间}\label{tab:macbyte}
\begin{tabularx}{\linewidth}{>{\raggedright\arraybackslash}p{3.0cm}>{\raggedright\arraybackslash}p{3.6cm}L}
\toprule
字节时间 & 线上内容 & 本拍动作 \\
\midrule
起始前一拍 & 空闲 & 判断起始，取样两网使能 \\
0--6 & \code{55} & \\
7 & \code{D5} & 取走第 0 个字节 \\
8 至 $8+L-1$ & 第 0 至 $L-1$ 个字节 & 取走下一字节，取走带 \code{s_last} 的字节后停止 \\
其后 4 拍 & FCS & A、B 两网各自更新 CRC32 \\
再后至少 12 拍 & 空闲 & 第 12 拍可判断下一帧起始 \\
\bottomrule
\end{tabularx}
\end{table}
